% Ban thu: chi giu cac cau co trong CD2

%@ id=TT25.11.003 | trang=1 | so=2.1 | de=II | vitri=CD2.4,C3.9.5 | muc=8 | tt=0 | soat=sua
\begin{bt}
Cho hai biểu thức $M=\dfrac{\sqrt{x}-2}{\sqrt{x}+1}$ và $N=\dfrac{\sqrt{x}}{\sqrt{x}+1}-\dfrac{\sqrt{x}-4}{1-x}$ với $x\ge 0$, $x\ne 1$.
\begin{cacy}(1)
\item Tính giá trị của biểu thức $M$ khi $x=8+2\sqrt{7}$.
\item Biết $P=\dfrac{M}{N}$. Chứng minh: $P=\dfrac{\sqrt{x}-1}{\sqrt{x}+2}$.
\item Đặt $P=\dfrac{M}{N}$. Tìm giá trị $x$ là số nguyên tố thỏa mãn $P\le \dfrac12$.
\end{cacy}
\loigiai{
\begin{cacy}(1)
\item Ta có: $x=8+2\sqrt{7}=(1+\sqrt{7})^2$ (thỏa mãn điều kiện).
Thay $x=(1+\sqrt{7})^2$ vào $M$. Tính được
\[M=\dfrac{1+\sqrt{7}-2}{1+\sqrt{7}+1}=\dfrac{\sqrt{7}-1}{\sqrt{7}+2}=\dfrac{(\sqrt{7}-1)(\sqrt{7}-2)}{(\sqrt{7}+2)(\sqrt{7}-2)}=3-\sqrt{7}.\]
\item Rút gọn
\[N=\dfrac{\sqrt{x}}{\sqrt{x}+1}-\dfrac{\sqrt{x}-4}{1-x}\]
\[N=\dfrac{\sqrt{x}(\sqrt{x}-1)}{x-1}+\dfrac{\sqrt{x}-4}{x-1}\]
\[N=\dfrac{x-\sqrt{x}+\sqrt{x}-4}{x-1}=\dfrac{x-4}{x-1}.\]
Suy ra
\[P=\dfrac{M}{N}=\dfrac{\sqrt{x}-2}{\sqrt{x}+1}:\dfrac{x-4}{x-1}=\dfrac{\sqrt{x}-1}{\sqrt{x}+2}.\]
Điều kiện xác định: $x\ge 0$, $x\ne 1$, $x\ne 4$.
\item Để $P\le\dfrac12$ thì $\dfrac{\sqrt{x}-1}{\sqrt{x}+2}\le\dfrac12$.
\[\dfrac{\sqrt{x}-1}{\sqrt{x}+2}-\dfrac12\le 0\]
\[\dfrac{\sqrt{x}-4}{\sqrt{x}+2}\le 0\]
$\sqrt{x}-4\le 0$ (vì $\sqrt{x}+2>0$ với mọi $x$ thuộc điều kiện xác định).
\[0\le x\le 16\]
Kết hợp điều kiện xác định, có: $0\le x\le 16$, $x\ne 1$, $x\ne 4$, mà $x$ là số nguyên tố nên suy ra $x\in\{2;3;5;7;11;13\}$.
\end{cacy}
}
\end{bt}
